\section{Implementation}

\subsection{Dataset Procurement}

\subsubsection{Object Detection Dataset}

The training dataset for the object detection model was sourced from the
Microsoft COCO dataset \citep{lin2014microsoft}. As mentioned prior, the COCO
dataset is a large-scale source of over $200,000$ labelled images across 80
object categories. The entirety of the dataset would lead to rundandancy with a
majority of classes being out of scope for the target application. Therefore, a
subset of the COCO dataset was created, containing only images including the
\texttt{person}, \texttt{backpack}, \texttt{handbag} and \texttt{suitcase}
classes, which I then aggregated into a single \texttt{luggage} class. This
resulted in a training dataset of \textbf{TODO} images split into training and
validation sets.

Originally, I had intended to use datasets specific to the abandoned luggage
detection domain, such as ABODA, however these datasets failed to include the
amount of variation needed for training a robust model. Many annotated datasets
created from the source footage, lead to overfitting and poor generalisation on
test examples. The COCO dataset, while lacking the domain specificity, offers a
much larger variety of examples, covering a wide range of environments,
lighting conditions, and object appearances.

Another benefit of using the COCO dataset is the ease of access and
availability of annotations that are compatible with the latest YOLO models. A
python script was used to filter the COCO dataset, specify the YOLO annotation
format, balance class instances, and oversample specific classes
% (\ref{sec:model_training_optimisation}) 
to create a more varied and proportionate training dataset.

\subsubsection{Suspicious Luggage Dataset}

% TODO: talk about the datasets to be used to test the suspicious luggage alert generation